\begin{annexitem}{annex:p15-goodman-short-partition}{thm:p15-goodman-short-partition}{Partition and nonconflict for the printed short condition}
\begin{proof}
Suppose $\operatorname{Ov}_\tau(J,H)$. Its final conjunct says there
is no eligible $L$ conflicting with $J$ and better entrenched than $J$.
If $J$ were overridden, the definition of $\operatorname{O}_\tau(J)$
would supply just such an $L$, together with further conditions.
Therefore $\neg\operatorname{O}_\tau(J)$. If $H$ were projectible,
its conflict with $J$ would require $\operatorname{O}_\tau(J)$,
contradicting this conclusion. Thus an overridden hypothesis is not
projectible. It is not nonprojectible either, since that definition
expressly requires a nonoverridden hypothesis.

For an eligible, nonoverridden $H$, classical logic gives two exhaustive
possibilities: every eligible conflict is overridden, or some eligible
conflict is not overridden. The former is exactly the projectible clause
and the latter the nonprojectible clause. The universal and existential
conditions cannot both hold. The overridden eligible hypotheses are the
remaining eligible unprojectible hypotheses. Every ineligible hypothesis
is unprojectible and satisfies neither definition requiring eligibility.
Hence the three classes cover $\mathcal H$ and are pairwise disjoint,
including when an eligible hypothesis has no eligible conflict.

If projectible $H,K$ conflict, projectibility of $H$ requires $K$ to be
overridden. The preceding argument shows that such a $K$ cannot be
projectible, a contradiction. No finiteness, numerical entrenchment scale
or normative adequacy premise was used.
\end{proof}
\end{annexitem}

\begin{annexitem}{annex:p15-goodman-parity-hierarchy}{thm:p15-goodman-parity-hierarchy}{The isolated path and the short-condition countercase}
\begin{proof}
Any upward hierarchy beginning at $H_i$ must next take $H_{i+1}$,
if that member exists: no other higher member conflicts with $H_i$.
Repeating the argument, maximality forces the hierarchy to be
$H_i,H_{i+1},\ldots,H_{n-1}$. It cannot extend upward beyond
$H_{n-1}$ since there is no external conflict. Each member conflicts
only with its adjacent members. It is therefore the unique maximal,
minimal hierarchy beginning at $H_i$, and it has $n-i$ members.
The source's parity condition gives the asserted overridden status.

If $n-i$ is odd, $H_i$ is not overridden and each existing neighbour
has an even-length suffix and is overridden. The projectible clause
holds. If $n-i$ is even, $H_i$ is overridden and so unprojectible; its
higher neighbour $H_{i+1}$ has an odd-length suffix and is not overridden,
so $H_i$ does not satisfy the projectible universal clause. Every
nonoverridden member has only overridden neighbours, leaving none
nonprojectible. For $n=1$ the sole member is nonoverridden and without
conflicts, hence projectible; the same reasoning covers that endpoint.

For $n=4$ under the short condition, $H_3$ is the only member without
a still-better conflicting rival. It overrides $H_2$. The member
$H_2$ cannot override $H_1$ because it conflicts with the better $H_3$;
$H_1$ cannot override $H_0$ because it conflicts with the better $H_2$.
Thus $H_2$ alone is overridden. Its only higher neighbour $H_3$ has
only an overridden conflict, making $H_3$ projectible. The members
$H_0,H_1$ are nonoverridden and conflict with one another, so both
are nonprojectible. The parity variant instead overrides $H_0,H_2$
and makes $H_1,H_3$ projectible. The constructions therefore differ
on this finite component. An arbitrary branching component is not
covered by the isolated-path proof.
\end{proof}
\end{annexitem}

\begin{annexitem}{annex:p15-goodman-regress-boundary}{thm:p15-goodman-regress-boundary}{A finite dependency bound and an infinite countercase}
\begin{proof}
For a node $v$, let $r(v)$ be the maximum length of an upward path
starting there. The possible lengths form a nonempty bounded set of
natural numbers, since the zero-length path exists and the finite-bound
hypothesis applies to $v$. Its maximum therefore exists. If $w$ is
directly above $v$, any upward path from $w$ can be preceded by the
step from $v$, increasing its length by one. Consequently $r(w)<r(v)$.

Induct on $r(v)$. A rank-zero node has no directly higher node, so it
cannot meet the necessary condition for final appreciability. For
positive rank, every directly higher node has lower rank and is not
finally appreciable by induction. Again the necessary condition
prevents $v$ from being finally appreciable. This proves the assertion
at each node, without needing a single uniform bound for the whole domain.

For the countercase use $0,1,2,\ldots$, with $i+1$ directly above
$i$. Mark none initially appreciable but mark all finally appreciable.
Each finally appreciable $i$ has the finally appreciable $i+1$ directly
above it. Thus the stipulated necessary condition holds at every node,
but the asserted absence of final appreciability is false. The paths
are unbounded, so the finite-bound hypothesis fails. A well-foundedness
condition or an otherwise justified additional interpretation is needed
to exclude this infinite self-sustaining assignment; it is not excluded
by the necessary condition alone.
\end{proof}
\end{annexitem}
